% !PRINTED 542
\weberpage{542}
If then $k$, $k_0$, $k'$ are the classes to which $\psi$, $\psi_0$, $\psi'$ belong, then
\[
\chi(\delta, k') = \chi(\delta, k_0)\,\chi(\delta, k),
\]
hence
\weq{5}{\begin{aligned}
\chi(\delta, k) &= (\delta, A), \\
\chi(\delta, k') &= (\delta, 2)(\delta, A).
\end{aligned}}

If $\omega_1$ is a root of the class $k$, then $\omega_2$ is a root of the opposite class $k^{-1}$; but in these two classes the characters $\chi(\delta, k)$ and $\chi(\delta, k^{-1})$ are the same. If we therefore multiply the formula \S\,141, (30) by $\chi(\delta, k)$ and form the sum over all classes $k$, it follows that:
\weq{6}{\overset{k}{\Sigma}\,\chi(\delta,k)\Bigl(\overset{x,y}{\Sigma}\frac{1}{\psi^s} - \overset{x,y}{\Sigma}\frac{1}{\psi'^s}\Bigr) = \frac{4\pi}{\sqrt{m}}\overset{k}{\Sigma}\,\chi(\delta,k)\log\frac{f(\omega)}{\sqrt[4]{2}},}
and on account of (5) one may write for the left-hand side:
% [sic] The printed fraction chi(delta,k)/1 is retained; its denominator is redundant, not an algebraic change.
\[
[1-(\delta,2)]\overset{k}{\Sigma}\frac{\chi(\delta,k)}{1}\overset{x,y}{\Sigma}\frac{1}{\psi^s}.
\]
Accordingly we have:
\weq{7}{\begin{aligned}
(\delta,2) = +1\colon &\quad \frac{2\pi}{\sqrt{m}}\overset{k}{\Sigma}\,\chi(\delta,k)\log\frac{f(\omega)}{\sqrt[4]{2}} = 0, \\
(\delta,2) = -1\colon &\quad \frac{2\pi}{\sqrt{m}}\overset{k}{\Sigma}\,\chi(\delta,k)\log\frac{f(\omega)}{\sqrt[4]{2}} = \overset{k}{\Sigma}\,\chi(\delta,k)\overset{x,y}{\Sigma}\frac{1}{\psi^s},
\end{aligned}}
where on the right the limiting value for $s=1$ is to be taken.

In \S\,113 we proved the formula:
\weq{8}{\overset{k}{\Sigma}\,\chi(\delta,k)\overset{x,y}{\Sigma}\frac{1}{\psi^s} = \Sigma\frac{(\delta,n)}{n^s}\Sigma\frac{(\delta',n)}{n^s}.}

If we denote by $K(\delta)$ the class number for the discriminant $\delta$, then, as was proved in \S\,112, (8): % faint pre-delta speck is scan/print damage, not a character
\weq{9}{\begin{aligned}
\sum^n \frac{(\delta,n)}{n} &= \frac{\pi}{\sqrt{-\delta}}\,K(\delta) \qquad \delta<0, \\
\sum^n \frac{(\delta,n)}{n} &= \frac{\log\varepsilon}{\sqrt{\delta}}\,K(\delta) \qquad \delta>0,
\end{aligned}}
in which the square roots are to be taken positive, and
\[
\varepsilon = \frac{T+U\sqrt{\delta}}{2}
\]
is the fundamental positive unit of the quadratic field with the fundamental number $\delta$. For the two exceptional cases $\delta=-3$, $\delta=-4$ these formulae hold if by $K$ one understands not the
%CONT>>
